\section{From local improvements to a global complexity theorem}
\label{sec:global}

\subsection{A complete cost ledger}

The distinction between normal weight and its encoding is essential but
insufficient. For a hierarchy state $H_i$, let $s_i$ be the full encoded
state size: source charts, current pattern, bundle identifications,
coordinate bits, and all live compressed-word or interval descriptions.
Let $a_i$ count candidate transitions that are generated, including later
discarded candidates. A useful global accounting identity has the form
\[
 T(n)\le T_{\rm input}(n)+
 \sum_{i\in\mathcal V}
 \left(T_{\rm state}(s_i)+
       \sum_{j=1}^{a_i}T_{\rm transition}(s_i,j)\right),
\]
where $\mathcal V$ contains every visited search state, including failed
branches. A polynomial bound on one summand does not control either
$|\mathcal V|$ or the maximum $s_i$.

The orbit-compression framework of Agol--Hass--Thurston supplies the right
kind of primitive for avoiding expansion in multiplicity
\cite{aht2006}. Lackenby's surface-curve algorithms likewise illustrate
operations charged to triangulation size and logarithmic curve weights
\cite{lackenby-curves}. The missing global work is to prove that the
particular hierarchy's interfaces stay within those input-size regimes
and that sufficiently few states and candidates are inspected. An
exponential number of short compressed objects is still too many objects.

\begin{proposition}[A sufficient search-tree accounting theorem]
\label{prop:global-ledger}
Suppose a deterministic recognition procedure on an $n$-crossing input
satisfies all of the following for some constants $a,b,c$:
\begin{enumerate}
\item The input conversion takes polynomial time.
\item Every visited state and every required transition certificate has
encoded size at most $n^a$.
\item At most $n^b$ children are generated at any state, and all local work
per child and per state is at most $n^c$.
\item Every root-to-leaf path in the actual explored tree has length at
most $d(n)$.
\item The branching rule is complete, the leaf predicates are sound, and
their combination correctly decides both knot answers.
\end{enumerate}
Then the total running time is $n^{O(d(n)+1)}$. In particular,
$d(n)=O(\log n)$ gives $n^{O(\log n)}$, while
$d(n)=O((\log n)^2)$ gives $2^{O((\log n)^3)}$.
\end{proposition}

\begin{proof}
The number of nodes is at most
$1+n^b+\cdots+(n^b)^{d(n)}$. Multiplying by a polynomial upper bound on
local work, including child generation and verification, gives the stated
bound. The final assertions follow by substituting the two depth estimates
and taking logarithms. The correctness hypothesis is listed separately
because a fast incomplete search would not prove a recognition theorem.
\end{proof}

The proposition is an accounting tool, not a claim that the hypotheses
already hold. In particular, it cannot be applied with the depth of one
successful path when the algorithm may search much deeper failed paths.
Memoization can improve this bound only after an authenticated equivalence
relation on states and a bound on distinct reachable equivalence classes
are proved.

\subsection{Strict quadratic decrease is weaker than multiplicative decrease}

Write $q=c_q(H)$ for Lackenby's quadratic complexity. The homological
certificate gives a strict raw drop of at least two under its stated
hypotheses. Along a retained forward chain whose other operations are
independently $q$-nonincreasing, there can therefore be at most
$\lfloor q_0/2\rfloor$ such strict cuts. This bound is useful: it excludes
indefinite repetition of these progress steps. It is generally linear in
the initial potential and gives neither logarithmic depth nor a bound on
the number of discarded attempts.

An attempted proof must also distinguish no progress from no available
certificate. The homology quotient may vanish, or a candidate may lie
inside the bundle image, while a geometric cut still lowers $q$. Conversely,
a nonzero quotient does not itself produce an admissible, efficiently
encoded connected normal surface. The quotient test supplies a sufficient
progress predicate after geometry has been authenticated.

There is a precise stronger target for a surface constructor. For a raw
cut, write $I_H\ge0$ for each parent zero-handle's integral index and
$i_{Ha}\ge0$ for the indices of all its children, so that
$\sum_a i_{Ha}=I_H$. Then $q=\sum_H I_H^2$.

\begin{proposition}[Balanced index splitting forces multiplicative progress]
\label{prop:balanced-index}
Suppose $q>0$. Let $0<\alpha\le1$ and $0<\delta<1$. Assume that a set
$\mathcal U$ of parent zero-handles satisfies
\[
 \sum_{H\in\mathcal U}I_H^2\ge\alpha q,
 \qquad
 \max_a i_{Ha}\le(1-\delta)I_H
 \quad(H\in\mathcal U).
\]
Then
\[
 q(H\cut S)\le(1-\alpha\delta)q(H).
\]
The same upper bound holds after any cleanup whose nonincrease of $q$
has been separately established.
\end{proposition}

\begin{proof}
For $H\in\mathcal U$, nonnegativity and index conservation give
\[
 \sum_a i_{Ha}^2\le (\max_a i_{Ha})\sum_a i_{Ha}
 \le(1-\delta)I_H^2.
\]
For any other parent, the same argument without the balance assumption
gives $\sum_a i_{Ha}^2\le I_H^2$. Summing proves a drop of at least
$\delta\sum_{H\in\mathcal U}I_H^2\ge\alpha\delta q$. The cleanup
assertion is immediate from its separately supplied nonincrease.
\end{proof}

For fixed positive $\alpha,\delta$, a forward nonincreasing chain can have
only $O(\log(q_0+1))$ such balance stages while its potential remains
positive. More precisely, after $k$ of them the potential is at most
$(1-\alpha\delta)^kq_0$; integrality bounds $k$ until the potential
vanishes or balance ceases to be available. This does not assert existence
of a next balance stage. An index-one parent cannot be split into
nonnegative integer children with maximum at most $(1-\delta)I_H$.
The residual small-index situation requires a separate geometric treatment,
and neutral stages and failed searches need their own bounds.

The inequality identifies the missing construction theorem more sharply
than the phrase ``find a useful surface.'' One would need to find a
compressed surface or a controlled batch that separates a fixed fraction
of the squared-index mass, or else returns a certified alternative region
whose handling admits an equally strong amortized estimate. An
outside-image witness guarantees some separation, but gives no uniform
fraction of the mass.

\subsection{A sharper raw-cut accounting potential}

The same identity gives a useful local reformulation. Define
\[
 \ell(H)=\sum_H I_H,\qquad
 P(H)=\sum_H\binom{I_H}{2}=\frac{q(H)-\ell(H)}{2}.
\]
Here $P$ counts unordered pairs of index units lying in the same parent.
This interpretation is arithmetic and does not assume those units have
canonical geometric representatives.

\begin{proposition}[Pair-separation accounting]
\label{prop:pair-potential}
For a raw normal cut satisfying nonnegative integral child indices and
per-parent index conservation,
\[
 \ell(H\cut S)=\ell(H),\qquad
 P(H)-P(H\cut S)=\sum_H\sum_{a<b}i_{Ha}i_{Hb}.
\]
Thus an outside-image cut from the progress theorem lowers $P$ by at least
one. Along a sequence consisting only of such raw cuts, the number of
strict cuts is at most $P(H_0)$.
\end{proposition}

\begin{proof}
Sum the conservation identities for the first equality. Subtract this
equality from the quadratic defect identity and divide by two for the
second. Each summand is a nonnegative integer. A positive quadratic defect
therefore gives an integral drop of at least one, and the final assertion
follows by telescoping the nonnegative potential.
\end{proof}

The qualification ``raw'' is necessary. Nonincrease of $q$ alone does not
imply nonincrease of $P$ if a subsequent operation changes $\ell$. For
example, as abstract index inventories, four indices equal to one and a
single index equal to two both have $q=4$, but their values of $P$ are zero
and one. This arithmetic example is not asserted to be an available
geometric cleanup. It shows why a proof of a cleanup rule for $q$ cannot
simply be relabelled as a proof for $P$.

The terminal condition $P=0$ means that every remaining index is zero or
one. It does not mean that the manifold is a ball, that its pattern is
essential, or that a knot is unknotted. A useful future strategy could
separate the problem of reducing index concentration from the classification
of this small-index remainder, provided both stages preserve the required
boundary data.

There is an even sharper count for a raw-only sequence. Let $r(H)$ be the
number of positive-index zero-handles. Every positive parent has at least
one positive child, every zero parent has only zero children, and a strict
raw cut has some parent with at least two positive children. Thus each
strict cut increases $r$ by at least one. Since $r\le\ell$ and $\ell$ is
conserved, the number of strict raw cuts is at most
\[
 \ell(H_0)-r(H_0)=\sum_H\max\{I_H-1,0\}\le P(H_0).
\]
The last inequality follows from $k-1\le\binom{k}{2}$ for every integer
$k\ge1$. The pair potential supplies the exact weighted defect; the
positive-child count supplies the sharper elementary number-of-cuts bound.
Neither extends through arbitrary cleanup without additional conservation
or monotonicity information.

\subsection{What the dual-forest theorem contributes to this ledger}

The top reduction resolves a specific local obstruction with polynomial
encoded work and a coordinate-magnitude-independent additional trace. It
does not add a large branching factor: the default uses one optimized tree.
Its optional complete restricted search has cost
$t^{O(v_3)}\operatorname{poly}(t,L)$ for $v_3$
interior source vertices, and hence is quasi-polynomial in the full encoded
source input when $v_3=O(\log t)$. This is a correct parameterized result for the stated
compiler family. No argument in this article makes $v_3$ logarithmic along
an arbitrary full recognition hierarchy or proves that that compiler family
is topologically complete.

Removing source three-handles also helps align a normal cut with interfaces
that require none. However, it preserves the specified bundle ports, not
an arbitrary hierarchy boundary pattern, and it does not automatically
establish Lackenby-admissibility. These remaining conditions are precisely
where a tempting informal composition would otherwise become invalid.
